\documentclass[11pt]{article}

% -----------------------------
% Packages (minimal, common)
% -----------------------------
\usepackage[a4paper,margin=1in]{geometry}
\usepackage{microtype}
\usepackage{hyperref}
\usepackage{enumitem}
\usepackage{graphicx}
\usepackage{booktabs}
\usepackage{amsmath,amssymb}

% -----------------------------
% Metadata
% -----------------------------
\title{STREAMTRACK: Multi-Camera Object Tracking Using Streams and Tables}
\author{Mahmoud Tayee, Ahmed Awad}
\date{\today}

\begin{document}
\maketitle

\begin{abstract}
% TODO: Replace with final abstract.
Multi-camera object tracking (MCOT) is typically implemented as a tightly coupled computer-vision pipeline, making it difficult to scale cross-camera association, to inspect intermediate decisions, and to adapt to changes in the camera network. This paper presents an alternative systems-oriented formulation: MCOT as a streaming program in which detections and tracklet updates are streams, identity hypotheses are maintained as relational state tables, and cross-camera association is expressed as continuous SQL/EPL queries. Inter-camera tracking is realised via topology-aware stream joins parameterised by a dynamically updateable camera neighbourhood table, reducing candidate explosion and enabling runtime reconfiguration without restarting the pipeline. 
\end{abstract}

\section{Introduction}
\subsection{Motivation}
% TODO: Add motivating use cases (smart buildings, retail, airports, traffic), and why multi-camera tracking is hard operationally.
Multi-camera object tracking must simultaneously satisfy (i) computer-vision accuracy requirements and (ii) operational requirements such as scalability, debuggability, and reconfigurability. Conventional implementations often embed cross-camera association logic into bespoke code, which complicates maintenance and makes runtime changes (e.g., adding a camera, changing neighbourhoods) disruptive.

\subsection{Key Insight}
% TODO: Make this crisp and aligned with your BEST streams+tables story.
We recast MCOT as a \emph{streams-and-tables} program: (a) \emph{streams} represent observations and incremental hypotheses; (b) \emph{tables} represent state (tracklets, global identities, topology/configuration); and (c) continuous queries implement the operator graph that transforms observations into persistent identity assignments.

\subsection{Contributions}
% TODO: Finalize claims once results are in.
\begin{itemize}[leftmargin=*]
    \item An operator-graph formulation of MCOT using streams and relational state tables.
    \item An end-to-end SQL/EPL realization over a CEP/streaming engine (Esper), including streaming counterparts of key clustering/association components.
    \item Topology-aware cross-camera association expressed as stream joins, reducing candidate comparisons relative to naive all-pairs matching.
    \item A dynamically updateable camera neighborhood topology table enabling runtime reconfiguration without code changes or pipeline resets.
    \item A performance-optimized implementation (\textbf{Ours}) that achieves a \textbf{25.6$\times$ speedup in similarity matrix generation} and a \textbf{4.3$\times$ overall speedup} in MCPT processing compared to the conventional baseline (\textbf{Theirs}).
\end{itemize}

\subsection{Paper Organization}
% TODO: Update after finalizing sections.
Section~\ref{sec:background} defines the problem and background. Section~\ref{sec:overview} presents the system overview and data model. Sections~\ref{sec:operators}--\ref{sec:joins} detail the operator graph and topology-aware joins. Section~\ref{sec:implementation} outlines the SQL/EPL implementation. Section~\ref{sec:evaluation} describes the evaluation plan. Sections~\ref{sec:discussion}--\ref{sec:conclusion} conclude with discussion, related work, and future directions.

\section{Background and Problem Definition}\label{sec:background}
\subsection{Multi-Camera Object Tracking (MCOT)}
% TODO: Formalize the inputs and outputs.
We consider a set of cameras $\mathcal{C} = \{c_1,\dots,c_m\}$. Each camera produces a stream of detections with associated appearance embeddings and optional metadata. The goal is to assign a \emph{global identity} to observations across cameras and time, producing consistent identity tracks.

\subsection{Baseline Pipeline (Starting Point)}
% TODO: Summarize the starting paper at high-level, mapping stages to your operators (no need for deep detail here).
We begin from a representative multi-camera tracking pipeline that (i) constructs tracklets per camera, (ii) selects representative images/observations, (iii) performs cross-camera re-identification (Re-ID) via clustering, and (iv) assigns lower-confidence tracklets in a subsequent stage. Our work keeps the functional intent but re-architects the pipeline as a streaming operator graph with explicit state tables and join-based cross-camera association.

\subsection{Streams-and-Tables Execution Model}
% TODO: Define only what is needed (streams, tables, windows, continuous queries).
A \emph{stream} is an unbounded sequence of tuples. A \emph{table} (relation) is a mutable stateful dataset updated incrementally by the stream. A \emph{continuous query} consumes streams, reads/updates tables, and emits derived streams or table updates.

% \subsection{Time Semantics: Ingestion Time}
% \label{subsec:ingestiontime}
% \textbf{We adopt ingestion time (processing-time) semantics.}
% All operators interpret time using the timestamp at which events are \emph{ingested by the engine}, not the original event-generation time. Windowing, ordering, and join constraints are defined on ingestion time. This simplifies runtime behavior under real-time video analytics deployments and avoids the need for watermarking and late-event handling.

% % TODO: If needed, later justify tradeoffs (e.g., late arrivals, jitter, vs operational simplicity).

\subsection{System Requirements}
% TODO: Convert into explicit requirements; later connect them to experiments.
\begin{itemize}[leftmargin=*]
    \item \textbf{Accuracy:} maintain competitive tracking and Re-ID quality.
    \item \textbf{Scalability:} handle increasing numbers of cameras and detections.
    \item \textbf{Reconfigurability:} change camera neighborhood topology without restarting/resetting.
    \item \textbf{Observability:} make intermediate state and decisions queryable for debugging.
\end{itemize}

\section{System Overview}\label{sec:overview}
\subsection{Architecture at a Glance}
% TODO: Add a figure of the operator DAG.
Figure~\ref{fig:dag} (placeholder) will depict the operator graph from detection ingestion to global identity updates.

\begin{figure}[t]
    \centering
    % TODO: Replace with actual diagram.
    \fbox{\parbox{0.9\linewidth}{\centering Placeholder: Operator DAG figure (streams as edges, tables as state).}}
    \caption{High-level operator graph for streams-and-tables multi-camera tracking.}
    \label{fig:dag}
\end{figure}

\subsection{Data Model: Streams and Tables}
% TODO: Replace with final schemas; keep names consistent with EPL/SQL later.
\paragraph{Streams (examples).}
\begin{itemize}[leftmargin=*]
    \item $\texttt{Detections}(\textit{camId}, \textit{ts\_ingest}, \textit{bbox}, \textit{score}, \textit{emb}, \textit{frameId}, \dots)$
    \item $\texttt{TrackletEvents}(\textit{camId}, \textit{trackletId}, \textit{ts\_ingest}, \textit{eventType}, \dots)$
    \item $\texttt{Representatives}(\textit{camId}, \textit{trackletId}, \textit{ts\_ingest}, \textit{repEmb}, \textit{quality}, \dots)$
    \item $\texttt{CrossCamCandidates}(\textit{srcCam}, \textit{dstCam}, \textit{srcTracklet}, \textit{dstTracklet}, \textit{ts\_ingest}, \textit{sim}, \dots)$
    \item $\texttt{IdentityUpdates}(\textit{globalId}, \textit{trackletId}, \textit{action}, \textit{confidence}, \textit{ts\_ingest}, \dots)$
\end{itemize}

\paragraph{State Tables (examples).}
\begin{itemize}[leftmargin=*]
    \item $\texttt{ActiveTracklets}(\textit{camId}, \textit{trackletId}, \textit{lastTs\_ingest}, \textit{state}, \textit{embStats}, \dots)$
    \item $\texttt{GlobalIdentities}(\textit{globalId}, \textit{members}, \textit{centroidEmb}, \textit{lastTs\_ingest}, \dots)$
    \item $\texttt{CameraTopology}(\textit{srcCam}, \textit{dstCam}, \textit{enabled}, \textit{delayMin}, \textit{delayMax}, \textit{gateParams}, \dots)$
    \item $\texttt{Calibration}(\textit{camId}, \textit{params}, \dots)$ (optional)
\end{itemize}

\subsection{Execution Semantics}
% TODO: Specify operator update triggers, state eviction policy, and window definitions in ingestion time.
Operators are evaluated continuously as new tuples arrive. Windows and joins use ingestion timestamps $\textit{ts\_ingest}$. State tables maintain current hypotheses and may apply eviction rules based on ingestion-time staleness (e.g., tracklet inactivity).

\section{Stream Operator Graph for MCOT}\label{sec:operators}
\subsection{Ingestion and Feature Enrichment Operators}
% TODO: Detail how detections are ingested and enriched with embeddings/pose/world coordinates.
This stage ingests per-camera detections and attaches feature vectors (appearance embeddings) and optional metadata (pose confidence, projected coordinates) as part of the streaming tuple.

\subsection{Single-Camera Tracking as Stateful Stream Processing}
% TODO: Formalize tracklet creation/update/end; specify association rules and state updates.
Single-camera tracking maintains $\texttt{ActiveTracklets}$ as a state table updated by detection streams. Output events include tracklet lifecycle events and incremental summaries.

\subsection{Streaming Conflict Resolution / Overlap Suppression (Optional)}
% TODO: Provide streaming equivalent of overlap suppression and how it is implemented with windowed constraints.
When detections/tracklets violate spatial overlap constraints within short ingestion-time windows, a conflict-resolution operator suppresses inconsistent hypotheses and triggers re-association.

\subsection{Representative Selection Operator}
% TODO: Define quality scoring and windowed ranking to pick representatives.
A continuous query selects representative observations for each tracklet based on quality criteria (e.g., embedding stability, pose confidence). The output is the $\texttt{Representatives}$ stream.

\subsection{Online Clustering / Global Identity Maintenance}
% TODO: Describe streaming clustering or incremental linkage approach.
Global identities are maintained in $\texttt{GlobalIdentities}$ via incremental updates (e.g., centroid maintenance, micro-clusters, merge rules). Identity updates are emitted as $\texttt{IdentityUpdates}$.

\section{Cross-Camera Tracking via Topology-Aware Stream Joins}\label{sec:joins}
\subsection{Naive Cross-Camera Matching Cost}
% TODO: Define baseline candidate generation and its complexity.
Naive cross-camera association compares tracklets across all camera pairs, leading to candidate explosion as $|\mathcal{C}|$ grows.

\subsection{Topology-Aware Candidate Generation}
% TODO: Provide join formulation and gating predicates (time, similarity, optional space).
We generate cross-camera candidates via a join constrained by the $\texttt{CameraTopology}$ table. Only camera pairs declared as neighbors are considered, and additional predicates restrict matches within ingestion-time delay windows and similarity thresholds.

\subsection{Camera Topology as a Reconfigurable State Table}
% TODO: Explain update mechanism and runtime effect.
The $\texttt{CameraTopology}$ table is treated as configuration state that can be updated online (e.g., enabling/disabling edges, modifying thresholds). Topology changes take effect immediately for subsequent ingested tuples, without code changes or pipeline restart.

\subsection{Join Planning and Optimization}
% TODO: Discuss selectivity, indexes/partitioning, and how topology sparsity helps.
Topology sparsity reduces the join search space and improves throughput. We later quantify candidate reduction and runtime cost savings relative to naive matching.

\section{Implementation in SQL/EPL on Esper}\label{sec:implementation}
\subsection{Why SQL-first on a CEP Engine}
% TODO: Argue for inspectability, modularity, and operational benefits.
A SQL/EPL-first design enables transparent intermediate states, rapid iteration, and deployment-friendly operations (monitoring, debugging, reconfiguration) compared to monolithic CV pipelines.

\subsection{Query Templates (Placeholders)}
% TODO: Add actual EPL snippets later. Keep structure stable.
\paragraph{(T1) Tracklet state maintenance.}
\begin{quote}
\emph{Placeholder for EPL/SQL that updates \texttt{ActiveTracklets} from \texttt{Detections} using ingestion-time windows.}
\end{quote}

\paragraph{(T2) Representative selection.}
\begin{quote}
\emph{Placeholder for EPL/SQL that selects top-$k$ representatives per tracklet in ingestion-time windows.}
\end{quote}

\paragraph{(T3) Topology-aware join.}
\begin{quote}
\emph{Placeholder for EPL/SQL that joins \texttt{Representatives} with \texttt{CameraTopology} and other \texttt{Representatives} to produce \texttt{CrossCamCandidates}.}
\end{quote}

\paragraph{(T4) Global identity updates.}
\begin{quote}
\emph{Placeholder for EPL/SQL that updates \texttt{GlobalIdentities} and emits \texttt{IdentityUpdates}.}
\end{quote}

\subsection{Streaming Counterparts of Clustering}
% TODO: Describe the chosen online clustering strategy and its implementation constraints in SQL/EPL.
We implement streaming counterparts of clustering by maintaining sufficient statistics (e.g., centroids, counts, similarity summaries) in tables and applying incremental merge/assignment rules.

\subsection{Handling Ordering and Jitter under Ingestion Time}
% TODO: Explain practicalities (buffers, small reorder windows) if needed.
Ingestion-time semantics avoids late-event complexity; minor jitter can be handled via small ingestion-time windows or buffering if required by the deployment.

\subsection{Configuration Parameters}
% TODO: List windows, thresholds, topology update rules.
We expose configurable parameters for window sizes, similarity thresholds, topology gating, and eviction policies, enabling controlled trade-offs between accuracy and performance.

\section{Experimental Evaluation}\label{sec:evaluation}
\subsection{Performance Benchmarking}
We benchmark the performance of the proposed streams-and-tables architecture (\textbf{Ours}) against the conventional baseline (\textbf{Theirs}). The evaluation is conducted on Scene 1 of the AIC24 Track 1 dataset using a window of 4 cameras and 192 candidate tracklets.

Table~\ref{tab:performance} summarizes the processing time for each stage of the Multi-Camera People Tracking (MCPT) pipeline. Our implementation significantly outperforms the baseline in similarity matrix generation, which is a major bottleneck in the original version.

\begin{table}[h]
    \centering
    \caption{Performance comparison of MCPT stages (\textbf{Theirs} vs. \textbf{Ours}).}
    \label{tab:performance}
    \begin{tabular}{@{}lrrr@{}}
        \toprule
        \textbf{Pipeline Stage} & \textbf{Theirs (ms)} & \textbf{Ours (ms)} & \textbf{Speedup} \\ \midrule
        Representative Selection & 230.70 & 221.15 & 1.04$\times$ \\
        Similarity Matrix Generation & 1260.00 & 49.27 & \textbf{25.57$\times$} \\
        Similarity Matrix Zeroing & 253.00 & 0.98 & \textbf{258.16$\times$} \\
        Similarity Replacement & $<$0.01 & $<$0.01 & 1.00$\times$ \\
        Hierarchical Clustering (HC) & 415.00 & 241.17 & \textbf{1.72$\times$} \\
        Global ID Assignment & 52.00 & 2.46 & \textbf{21.13$\times$} \\ \midrule
        \textbf{Total MCPT Time} & \textbf{2210.70} & \textbf{515.99} & \textbf{4.28$\times$} \\ \bottomrule
    \end{tabular}
\end{table}

\subsection{Performance Analysis}
The results highlight a significant breakthrough in \emph{Similarity Matrix Generation}, where \textbf{Ours} achieves a 25.57$\times$ speedup. This is attributed to the efficient handling of primitive arrays and JIT-optimized execution of cosine similarity loops. Furthermore, the Hierarchical Clustering (HC) stage now outperforms the baseline by 1.72$\times$ after optimizing for production mode by disabling deterministic sortings required for debugging. Overall, the system achieves a processing time of 516 milliseconds for a 4-camera window, demonstrating its robust capability for real-time operation in complex multi-camera environments.

\subsection{Research Questions}
\begin{itemize}[leftmargin=*]
    \item \textbf{RQ1:} Does the streams-and-tables re-architecture preserve tracking quality compared to the baseline pipeline?
    \item \textbf{RQ2:} How much does topology-aware joining reduce candidate comparisons and runtime cost compared to naive all-pairs matching?
    \item \textbf{RQ3:} How robust is the system under runtime topology changes (without restart/reset) in terms of stability and accuracy?
    \item \textbf{RQ4:} How do throughput and latency scale with the number of cameras and detection rates?
\end{itemize}
\subsection{Datasets and Scenarios}
% TODO: Fill with selected datasets and synthetic topology-change scenarios.
We will evaluate on (i) publicly available multi-camera tracking datasets and (ii) controlled scenarios that vary camera graph sparsity and topology updates during execution.

\subsection{Metrics}
% TODO: Choose final set (IDF1, HOTA, ID switches; systems metrics).
We report standard tracking metrics (e.g., IDF1/HOTA, identity switches) and systems metrics (throughput, end-to-end ingestion-time latency, memory footprint, join candidate counts).

\subsection{Baselines}
% TODO: Add baseline implementations.
Baselines include the starting pipeline, a naive cross-camera matcher without topology pruning, and ablations disabling key operators (representative selection, clustering variants).

\subsection{Ablation Studies}
% TODO: Specify planned ablations and expected outcomes.
We conduct ablations to isolate the effects of topology pruning, representative selection, and streaming clustering choices.

\section{Discussion}\label{sec:discussion}
\subsection{Benefits of a SQL/Streams-and-Tables Formulation}
% TODO: Discuss inspectability, modularity, runtime reconfiguration, observability.
We discuss how expressing MCOT as continuous queries improves modularity, debuggability, and operational control, particularly under changing camera networks.

\subsection{Limitations and Failure Modes}
% TODO: e.g., ambiguous Re-ID, topology misconfiguration, embedding drift.
We document limitations including challenging Re-ID conditions, sensitivity to topology configuration, and appearance drift, and discuss mitigation strategies.

\subsection{Deployment Considerations}
% TODO: Edge/cloud partitioning, monitoring, alerts, etc.
We outline deployment patterns and monitoring hooks, such as join selectivity dashboards, identity churn indicators, and topology update audit logs.

\section{Related Work}\label{sec:related}
% TODO: Add citations and organize by buckets.
We will cover (i) multi-camera tracking and Re-ID pipelines, (ii) stream processing for video analytics, (iii) online/incremental clustering, (iv) CEP engines and continuous query systems, and (v) topology-aware processing and constrained joins.

\section{Conclusion and Future Work}\label{sec:conclusion}
% TODO: Replace with final conclusion once results are available.
We presented a draft of a streams-and-tables re-architecture of multi-camera object tracking implemented in SQL/EPL on a CEP engine, featuring topology-aware stream joins and runtime reconfigurable camera neighborhoods under ingestion-time semantics.

\paragraph{Future work.}
% TODO: Expand as desired.
Future directions include multi-hop topology reasoning, learning topology weights, richer spatio-temporal constraints, and edge-deployed variants.

\bibliographystyle{plain}
% \bibliography{references} % TODO: add later

\end{document}
